\documentclass[11pt]{article}
\usepackage[margin=1in]{geometry}
\usepackage{amsmath,amssymb,amsthm,mathtools,hyperref}
\newtheorem{proposition}{Proposition}
\newtheorem{definition}{Definition}
\newtheorem{remark}{Remark}
\newcommand{\Prob}{\mathbb P}
\newcommand{\E}{\mathbb E}
\newcommand{\R}{\mathbb R}
\newcommand{\ind}{\mathbf 1}
\newcommand{\TV}{\operatorname{TV}}
\newcommand{\VaR}{\operatorname{VaR}}
\newcommand{\ES}{\operatorname{ES}}
\DeclareMathOperator*{\argmax}{arg\,max}
\title{Certificate Chains for Financial Text Models:\\Conditional Results and Corrections to the Initial Draft}
\author{130U-FinLLM research notes, WP0.3}
\date{15 November 2026}
\begin{document}
\maketitle
\begin{abstract}
This note formalizes the initial T0, T1.1--T1.5,
T2.1--T2.4 and T3.1. Statements are conditional mathematical results,
not claims of novelty or real-data certification. In particular, test
success is not a universal proof; recorded-time selection can leak future
information; the shard-count lower bound has a tie-breaking exception;
numeric monotonicity only applies to designated coordinates with fixed
context; and marginal calibration does not identify individual error
probabilities. The accompanying Chinese file \texttt{notes.md} contains
the full operational contract and decision implications. This source was
not compiled: no TeX engine was available in the inspected environment.
\end{abstract}

\section{T0: Semantics and composition}
A certificate records a claim $\varphi$, assumptions $A$, evidence $E$,
and a versioned checker $V$. Deterministic soundness is the universal
implication $V(E)=\mathrm{accept}\Rightarrow(A\Rightarrow\varphi)$.
Finite tests alone do not prove this implication. A statistical
certificate additionally records the sampling mechanism and a coverage
level. Coverage of a procedure need not equal coverage conditional on
selecting an accepted output.
\begin{proposition}[Certificate composition]
On a specified common probability space, suppose
$\Prob(E_i^c)\le\alpha_i$ and, under $A$, $\bigcap_i E_i$ implies
$\varphi$. Then $\Prob(\varphi)\ge1-\min(1,\sum_i\alpha_i)$.
\end{proposition}
\begin{proof}
The failure event is contained in $\bigcup_i E_i^c$; apply the union
bound. Disjoint failure events show that no uniformly better bound is
available without dependence information. Independence is unnecessary.
\end{proof}
Assigning zero failure probability to a mathematical implication does
not assign zero probability to implementation bugs or false premises.

\section{T1.1: History, randomness, and selection}
Treat content versions as atomic records. Let $\mathcal F_t$ contain
all actually available versions, IDs, collection metadata and selection
inputs at time $t$. If $H_t$ is a history snapshot and $\Omega$ is
training randomness, then a measurable map
$\theta=F(H_t,\Omega)$ is
$\mathcal F_t\vee\sigma(\Omega)$-measurable. For each fixed seed,
equal histories give equal outputs. Randomized parameters need not be
literally $\mathcal F_t$-measurable unless randomness is included in the
initial information set.

These measurability and fixed-seed statements do not themselves require
independence. To assert that the seed adds no future information, assume
it is independent of the entire data-generating world, or conditionally
independent of the future given $\mathcal F_t$. Marginal independence
from the future alone is insufficient: with independent fair bits $H,B$,
let $\Omega=H\mathbin{\mathrm{xor}}B$. Then $\Omega$ is independent of
$B$, yet $(H,\Omega)$ recovers the future bit $B$. The conditional
noninterference proof and domain restrictions are expanded in the
companion \texttt{pit\_noninterference.md} note (2026-11-16 review).

Information measurability is distinct from admissible historical provenance.
A text constant in the chosen probability space, or a realized checkpoint
subsequently treated as a constant, can be measurable without proving that
its bytes were available or its computation was admissible at the cutoff.
Historical construction therefore also needs an ex ante fixed algorithm,
cross-world data semantics, and version-availability and selection contracts;
it cannot be obtained by conditioning away the realized training process.

Write $D_{\rm obs}=\{i:\tau_i\le t\}$ and
$C_t=\{i:\tau_i\le t<\tau_i^*\}$. The deletion comparator uses
$D_{\rm del}=D_{\rm obs}\setminus C_t$ with the same algorithm and
randomness. It is historical only if its selection, metadata and shared
components are functions of $H_t$. It need not equal retraining on all
truly available historical data.

\paragraph{Counterexample to unconditional deletion sufficiency.}
Both $d_0,d_1$ exist before $t$. A future bit $B$ determines which one
receives recorded timestamp $t-1$; the other is assigned $t+1$.
Record-based filtering retains only $d_B$, and a learner outputs its
index. Here $C_t=\varnothing$, yet the output reveals $B$. Missing old
documents can therefore carry information if their selection uses the
future. A historical selection rule repairs this gap by measurable
composition. Tokenizers, deduplication, model selection, labels,
teachers, stopping rules and all auxiliary lookups belong in this rule.

\section{T1.2: DP expectation bounds}
Assume the complete algorithm $M$ is $(\varepsilon,\delta)$-DP under
whole-document addition or deletion. Let $D,D'$ be connected by at
most $k$ such operations, and fix the same measurable $S\in[0,1]$ for
both output laws. Set
\[
a=k\varepsilon,\qquad b=\delta\sum_{j=0}^{k-1}e^{j\varepsilon},
\]
with $b=0$ when $k=0$.
\begin{proposition}[Symmetric expectation envelope]
\[
|\E S(M(D))-\E S(M(D'))|
\le\min\left\{1,\frac{e^a-1+2b}{e^a+1}\right\}.
\]
\end{proposition}
\begin{proof}
Iterate the DP inequality along the adjacency path to obtain
$P(A)\le e^a Q(A)+b$ and its reverse. If
$d=P(A)-Q(A)\ge0$, $q=Q(A)$, apply the forward inequality to $A$ and
the reverse inequality to $A^c$:
\[
d\le(e^a-1)q+b,\qquad e^a d\le(e^a-1)(1-q)+b.
\]
Adding yields $(1+e^a)d\le e^a-1+2b$. Swap laws and take the
supremum over events. Layer-cake integration of
$\E S=\int_0^1\Prob(S>u)\,du$ gives the expectation bound.
\end{proof}
The weaker $\min(1,e^a-1+b)$ is also valid, with
$b\le k e^{(k-1)\varepsilon}\delta$ for $k\ge1$ and $b=k\delta$
when $\varepsilon=0$. These are bounds between randomized training
laws, not deterministic intervals for one realized checkpoint.

For $\mu$-GDP define
$G_\mu(\alpha)=\Phi(\Phi^{-1}(1-\alpha)-\mu)$.
Adjacent laws obey, for every event or randomized test,
$P_{j+1}(A)\le\Phi(\Phi^{-1}(P_j(A))+\mu)$.
Iteration gives the same bound with $k\mu$. Thus the minimum sum of
testing errors is at least $2\Phi(-k\mu/2)$ and
\[
|\Delta\E S|\le\TV\le2\Phi(k\mu/2)-1.
\]
Requiring this to be at most $.05$ gives
$k\mu\le2\Phi^{-1}(.525)\simeq .1254$.
No utility impossibility for a 24M model follows from this calculation.
Token-level privacy is not automatically document-level privacy.

\section{T1.3: Fixed partitions and statistic envelopes}
There are $m\ge1$ heads, $K\ge2$ labels and smallest-label tie breaking.
Shared representation, tokenizer, hyperparameters and features are fixed
from a clean historical core. A permanent ID-to-shard map is unchanged
between comparators. Each head is a deterministic function of its own
shard and exogenous randomness; there is no cross-shard fitting state.
Deleting $k$ records or modifying $k$ records at fixed IDs touches at
most $k$ shards. Replacing IDs can touch $2k$ shards instead.

\begin{proposition}[Vote certificate]
Let $n_c(z)$ count votes and $a$ be the observed winner. If, for all
$c\ne a$,
\[
n_a(z)-n_c(z)\ge2b+\ind\{c<a\},
\]
then the winner survives arbitrary changes of at most $b$ heads.
\end{proposition}
\begin{proof}
One changed vote reduces a pairwise margin by at most two. After $b$
changes the margin is still positive against lower-index opponents and
nonnegative against higher-index opponents, exactly as tie breaking requires.
\end{proof}

Fix evaluation points $z_j$ and scores $s_j(c)\in[0,1]$, $j=1,\ldots,J$.
Let $R_j(b)$ contain all labels reachable by at most $b$ vote edits;
a computational relaxation must be an outer set. Then
\[
L(b)=\frac1J\sum_j\min_{c\in R_j(b)}s_j(c),\qquad
U(b)=\frac1J\sum_j\max_{c\in R_j(b)}s_j(c)
\]
enclose the score of every common set of at most $b$ altered shards.
Indeed each realized label lies in its reachable set, and summation
preserves inequalities. Optimizing separately over points is a relaxation
of a common corrupted-shard set. Macro-F1, Sharpe ratios and drawdowns are
not this additive statistic and require separate envelopes.

\section{T1.4: Exact finite-population auditing}
There are $N$ fixed suspect records and $K^*$ truly contaminated records.
Before seeing audit outcomes, choose $0\le n\le N$ and draw a simple
random sample without replacement. Audits are assumed error-free. Let
$X=x$ be the count of contaminated sampled records and $0<\alpha<1$.
Define
\[
F_K(x)=\sum_{j=\max(0,n-N+K)}^{\min(x,n,K)}
\frac{\binom Kj\binom{N-K}{n-j}}{\binom Nn},\qquad
K_u(x)=\max\{0\le K\le N:F_K(x)>\alpha\}.
\]
An empty sum is zero; $K_u=N$ for $n=0$ and $K_u=x$ for $n=N$.

\begin{proposition}[Uniform coverage]
For every fixed $K^*$, $\Prob_{K^*}(K^*\le K_u(X))\ge1-\alpha$.
\end{proposition}
\begin{proof}
Couple populations by changing good records to bad ones, keeping the
same random sample. Then $X_K$ increases with $K$, so $F_K(x)$ decreases.
Failure $K_u(X)<K^*$ implies $F_{K^*}(X)\le\alpha$.
The event comprises an initial segment of the discrete support; its
probability is the CDF at its largest point, at most $\alpha$.
\end{proof}

For zero discoveries, when $K\le N-n$,
\[
F_K(0)=\frac{\binom{N-K}n}{\binom Nn}
=\prod_{i=0}^{n-1}\left(1-\frac K{N-i}\right)
\le(1-K/N)^n\le e^{-nK/N};
\]
otherwise $F_K(0)=0$. Consequently every accepted $K$ is strictly less
than $N(1-\alpha^{1/n})$ and $N\log(1/\alpha)/n$ when $n>0$.
The approximation $K_u\simeq N\log(1/\alpha)/n$ needs a small sampling
fraction as well as a small contamination fraction. It is the scale of
an upper confidence bound after zero discoveries, not a lower bound on
the number of bad records actually present. For $N=1000,n=500,\alpha=.05$
the exact integer bound is $4$, whereas the rule-of-three approximation
is about $6$. These are hypothetical population sizes, not measured SEC data.

\paragraph{Tie-breaking correction.}
The assertion that any robust prediction requires $m\ge2b+1$ is false:
with $m=2,b=1$, both votes for label 1, any one vote edit leaves label 1
winning or tied and preferred. In general $m=2b$ suffices for a unanimous
minimum-label prediction. The displayed margin condition requires
$m\ge2b$ if $a=1$, and $m\ge2b+1$ if $a>1$. Strict majority regardless
of priority requires the latter; actual margins must still be checked.

\paragraph{Fixed-budget audit gate.}
For an integer $0\le B<N$, CDF monotonicity gives exactly
\[
K_u(0)\le B\quad\Longleftrightarrow\quad
\frac{\binom{N-(B+1)}n}{\binom Nn}\le\alpha.
\]
Thus the smallest integer $n$ satisfying this inequality is the required
zero-discovery audit size. Binomial coefficients outside support are zero.
The margin budget at point $j$ is
\[
B_j=\left\lfloor\frac12\min_{c\ne a_j}
\{n_{a_j}-n_c-\ind\{c<a_j\}\}\right\rfloor.
\]
For a model frozen before auditing, certifying at least $r$ points is
ensured by using the $r$-th largest $B_j$ as the budget gate.

\paragraph{Prediction robustness and nonvacuous statistics differ.}
With votes $(1,1,2,2,3)$ and $b=1$, reachable winners are $\{1,2\}$.
The prediction is not robust, but if the true label is 3 its accuracy is
exactly zero. Thus an unrobust prediction can still contribute a point
interval rather than $[0,1]$. Collective constraints can also tighten
aggregate intervals. A valid interval is nonempty; ``nonvacuous'' means
narrower than the trivial range.

\paragraph{Stratification and optional stopping.}
For fixed source strata, obtain bounds $K_{u,h}$ at prescribed levels
$\alpha_h$ with $\sum_h\alpha_h\le\alpha$. The union bound yields
$K^*\le\sum_hK_{u,h}$ with confidence $1-\alpha$, without independence.
Unobserved strata have upper bound $N_h$.
Repeated use of a fixed-time interval after optional stopping is invalid
without correction. One valid zero-discovery process for $H_0:K\ge K_0$
is
\[
E_n=\ind\{X_n=0\}/F_{K_0,n}(0),\quad0\le n\le N-K_0.
\]
On a zero path the next conditional no-error probability under $K$ is
$(N-K-n)/(N-n)$, no greater than its $K_0$ counterpart. After a discovery
the process stays zero. Hence it is a nonnegative supermartingale starting
at 1. Stopping at its first crossing of $1/\alpha$ shows crossing probability
at most $\alpha$. Zero-path crossing is precisely $F_{K_0,n}(0)\le\alpha$;
optional stopping does not give a free improvement of that constant.

\section{T1.5: The deletion-comparator chain}
Assume: historical selection and all shared components (T1.1); a clean
shared core; a suspect universe frozen before audit; fixed-shard local
training (T1.3); valid error-free random auditing (T1.4); fixed additive
bounded evaluation scores and denominator; and sound computation of an
outer envelope, including any required floating-point error control.

\begin{proposition}[Conditional finite-test-set certificate]
With probability at least $1-\alpha$ over auditing,
$S(F_{\rm del})\in[L(K_u),U(K_u)]$.
\end{proposition}
\begin{proof}
On $K^*\le K_u$, deletion affects at most $K^*$ shards. Every unaffected
head sees identical inputs, shared parameters and randomness, so its
output is identical. T1.3 encloses the comparator score. T1.4 bounds the
probability of the exceptional audit event. The historical assumptions
are needed to interpret the comparator as look-ahead free.
\end{proof}

A single contaminated core record may change the shared representation
and all heads. Adding a core contamination upper bound to an edge count
does not preserve the shard argument. A core audit can instead establish
zero core contamination with additional failure budget, or a separate
shared-backbone stability theorem is needed.
If evaluation observations are additionally iid, independent of training
and audit, and the clean comparator is fixed before evaluation, Hoeffding
adds $\epsilon_J=\sqrt{\log(2/\beta)/(2J)}$ to each endpoint for a
population-mean interval with confidence $1-\alpha-\beta$. General cluster
bootstrap does not automatically supply an exact finite-sample substitute.

\section{T2.1: Behavioral relations}
For an edit $e$, define distribution preservation by
$\TV(q_e,q)\le\rho$; label preservation by equal tie-broken argmax;
relative relabeling margin by a change of at least $\gamma$ in the
$b$-versus-$a$ log odds (positive relevant probabilities required);
UP by nondecreasing upper-tail probabilities; and FREE by no constraint.
These are distinct specifications.
UP need not preserve labels: $(.6,.3,.1)\to(.2,.3,.5)$ is an upward
stochastic shift with a different argmax. Label preservation need not
give UP: $(.2,.6,.2)\to(.3,.6,.1)$. A relabeling margin does not force
a label switch: $(.99,.01)\to(.98,.02)$ increases relative odds but keeps
the winner. Distribution preservation implies label preservation if the
original top-two gap exceeds $2\rho$, because each coordinate changes by
at most $\rho$. Neutral strengthening can move probability inward from
both tails, and is not a stochastic upward shift.

\section{T2.2: Ordinal monotonicity and numeric sensitivity}
Fix context $c$, $g\in\R^d$, $M\subseteq\{1,\ldots,d\}$ and
$\lambda\ge0$. Suppose $h(\cdot,c)$ is $\lambda$-Lipschitz in $\ell_1$.
Set $\eta=h(g,c)+\lambda\sum_{r\in M}g_r$ and
\[
Q_i=\Prob(Y\le i\mid g,c)=\sigma(\theta_i-\eta),\qquad
\theta_1<\cdots<\theta_{K-1},
\]
with thresholds fixed under numeric interventions.

\begin{proposition}[Corrected monotone-head statement]
If $g'_r\ge g_r$ for $r\in M$ and $g'_r=g_r$ otherwise, the output
distribution at $g'$ first-order stochastically dominates that at $g$.
At differentiability points, for $r\in M$,
\[
0\le\partial_r\eta\le2\lambda,\quad
0\le\partial_rq_K\le\lambda/2,\quad
-\lambda/2\le\partial_rq_1\le0,\quad
|\partial_rq_i|\le\lambda/2\ (1<i<K).
\]
For arbitrary $g,g'$ with fixed context,
$|q_i(g')-q_i(g)|\le(\lambda/2)\|g'-g\|_1$ at every point.
Without the residual term the constant is $\lambda/4$, with no direction
guarantee.
\end{proposition}
\begin{proof}
For an allowed nonnegative increment $\Delta$,
$\Delta h\ge-\lambda\|\Delta\|_1$ while the residual increases by
$\lambda\|\Delta\|_1$. Thus $\eta$ increases, and all $Q_i$ decrease.
One-coordinate difference quotients bound $\partial_r\eta$ by $0$ and
$2\lambda$ wherever the derivative exists; Lipschitzness guarantees
almost-everywhere differentiability. Each endpoint class probability
has derivative in $\eta$ of magnitude at most $1/4$. For an intermediate
class the derivative is a difference of two numbers in $[0,1/4]$, so
its magnitude is also at most $1/4$. Apply the chain rule. Finally $\eta$
is $2\lambda$-Lipschitz and each class probability is $1/4$-Lipschitz
in $\eta$, proving the finite-increment statement everywhere.
\end{proof}

The restriction on nonmonotone coordinates is essential: take $M=\{1\}$,
$h(g)=-\lambda g_2$. Context independence is also essential: if
$c=g$ and $h(g,c)=-2\lambda c$, then $h$ is constant in $g$ for each
fixed $c$, but the actual $\eta=-\lambda g$ decreases. Numeric words,
units and parser side channels must therefore be covered by the contract.
If $g=(x-b)/b$ with fixed $b>0$, the raw-value derivative is bounded by
$\lambda/(2b)$; there is no uniform bound near $b=0$, and negative $b$
reverses direction. These structural results do not require the context
encoder to be trained from scratch.

\section{T2.3: Units and complete invariants}
Let $x\in\R_{>0}^p$ and let $D\in\R^{q\times p}$ be the dimension
matrix. Positive unit scales act by
$x'_j=x_j\prod_a s_a^{D_{aj}}$. For $v\in\ker D$,
$\Pi_v(x)=\prod_jx_j^{v_j}$ is invariant since the scale exponents
are $Dv=0$. A model depending only on such features is invariant
\emph{provided its context and all auxiliary paths are invariant too}.

For the expressivity claim, take a complete basis $B$ of $\ker D$.
If $B^T\log x=B^T\log y$, then
$\log y-\log x\in(\ker D)^\perp=\operatorname{im}D^T$.
Hence $x,y$ lie in the same unit orbit. Any orbit-invariant function
therefore factors through the complete invariants. Arbitrarily selecting
a few ratios need not separate orbits and can lose expressivity.
Zeros and signed financial quantities require separate zero/sign strata
and nonzero denominators; logarithms and arbitrary powers are not valid
on the full signed domain. Unit tokens must be normalized as well.
Real-number invariance is not itself a claim of bitwise floating-point
invariance. Currency market movements are not pure reporting-unit changes.

\section{T2.4: Entity slots}
For a specified renaming family $\Pi$, assume a canonicalizer $a$ satisfies
$a(\pi x)=a(x)$ for all $x,\pi$. Then $F=G\circ a$ obeys
$F(\pi x)=F(x)$ by substitution. This does not prove that aliases,
coreference, targets and slot order are parsed correctly. Naming slots
by alphabetic entity order can violate the premise. Finite alias tests
only establish facts about the tested alias family.

\section{T3.1: Bernoulli dependence uncertainty}
Within each finite cluster, error indicators have known true marginals
$e_i\sim\operatorname{Bernoulli}(p_i)$; different clusters are independent.
Let $e_i^{\rm com}=\ind\{U_c\le p_i\}$ within cluster $c$, using
independent uniform $U_c$ across clusters.
\begin{proposition}[Worst convex risk]
For $E=\sum_i e_i$ and $E^{\rm com}=\sum_i e_i^{\rm com}$,
$E\le_{\rm cx}E^{\rm com}$. Variance, stop-loss and upper-tail ES
are maximized by this coupling.
\end{proposition}
\begin{proof}
For one cluster represent a joint law by masses $q_A$ on subsets $A$
of erroneous coordinates. Marginal constraints define a compact
polytope. For any convex $\phi$, choose a law maximizing
$\sum_Aq_A\phi(|A|)$; among maximizers, maximize $\sum_Aq_A|A|^2$.
If two incomparable sets $A,B$ have positive mass, move
$\epsilon=\min(q_A,q_B)>0$ mass from each to $A\cap B,A\cup B$.
This preserves every marginal and total probability. Convexity ensures
the primary objective cannot decrease; the secondary objective increases
by $2\epsilon|A\setminus B||B\setminus A|>0$, a contradiction.
Thus optimal positive-mass sets form a chain, whose masses are the
successive differences of ordered $p_i$ and hence are the comonotone law.

If $X\le_{\rm cx}Y$ and $Z$ is independent, the map
$x\mapsto\E\phi(x+Z)$ is convex, proving $X+Z\le_{\rm cx}Y+Z$.
Replace clusters successively. The mean is fixed by the marginals, so
the $x^2$ inequality maximizes variance; positive-part functions yield
stop-loss. Finally, for $0\le\beta<1$,
\[
\ES_\beta(X)=\inf_t\left[t+\frac{\E(X-t)_+}{1-\beta}\right].
\]
Every objective at fixed $t$ is ordered, and infima preserve that order.
\end{proof}

Sorted $p_i$ together with 0 and 1 determine each cluster's exact
comonotone count distribution; convolution gives the independent-cluster
sum. VaR is different: for two marginals $p_i=.5$, the comonotone sum
has equal masses at 0 and 2, so under
$\VaR_\beta(X)=\inf\{x:F_X(x)\ge\beta\}$ its median VaR is 0.
The antimonotone sum is identically 1, with median VaR 1.

Calibration alone does not supply known $p_i$. Two equal-sized subgroups
with error rates 0 and .2 can both be assigned .1 and still be calibrated
in aggregate. Plug-in probabilities therefore give a conditional scenario,
not an unconditional risk certificate. Simultaneous upper bounds
$p_i\le u_i$ can bound increasing convex risks via comonotone coupling
at $u_i$, but not variance by that monotonicity argument. If clusters
are dependent, merging them is conservative for convex risks; independent
convolution is unjustified.

\section{Computational scope and open evidence}
The standard-library script \texttt{research/checks/theory\_checks.py}
checks finite voting cases, an exact rational linear-head contract,
100,000 floating-point points, unit rescalings, and a finite grid of
Bernoulli couplings. Its output records hashes and timing. It does not
train a transformer or certify an actual parser, data source or deployment.
The Bernoulli theorem covers all couplings by proof; the grid enumeration
does not. Real-data nonvacuity, prior-art coverage, parsing error and G0/G1
remain separate empirical or research obligations.

\begin{thebibliography}{9}
\bibitem{gdp} J. Dong, A. Roth, W. J. Su.
\emph{Gaussian Differential Privacy}.
\url{https://arxiv.org/abs/1905.02383}.
Metadata and abstract accessed; mathematical calculations are supplied here.
\bibitem{monotone} O. Kitouni, N. Nolte, M. Williams.
\emph{Robust and Provably Monotonic Networks}.
\url{https://arxiv.org/abs/2112.00038}.
Metadata and abstract accessed; no novelty claim is made for residual monotonicity.
\end{thebibliography}
\end{document}
